\documentclass{beamer}

\mode<presentation>





\usetheme{Frankfurt}%
\usecolortheme{seagull}
\logo{\includegraphics[height=.25in]{clarksonGreen}}

\definecolor{garnet}{RGB}{136,0,0}
%\definecolor{clarksonGreen}{RGB}{0,71,28}
\definecolor{clarksonGreen}{RGB}{0,52,21}
\setbeamercolor{palette primary}{fg=clarksonGreen,bg=white}
\setbeamercolor{palette secondary}{fg=clarksonGreen,bg=white}
\setbeamercolor{palette tertiary}{fg=clarksonGreen,bg=white}
\setbeamercolor{palette quaternary}{bg=clarksonGreen,fg=white}
\setbeamercolor{block title}{fg=black,bg=black!15}
\setbeamercolor{block body}{fg=black,bg=black!10}
\setbeamercolor{titlelike}{bg=clarksonGreen,fg=white} % parent=palette quaternary}

\newcommand{\half}{\mbox{$\frac{1}{2}$}}
\newcommand{\deltat}{\mbox{$\triangle t$}}
\newcommand{\deltax}{\mbox{$\triangle x$}}
\newcommand{\deltay}{\mbox{$\triangle y$}}

\newcommand{\deriv}[2]{\frac{d}{d#2}#1}
\newcommand{\derivTwo}[2]{\frac{d^2}{d#2^2}#1}

\newcommand{\lp}{\left(}
\newcommand{\rp}{\right)}


\newcommand{\arrayTwo}[4]{
  \left[
  \begin{array}{rr}
    #1 & #2 \\
    #3 & #4
  \end{array}
  \right]
}


\newcommand{\startRowOpsTwo}{
  \left[
    \begin{array}{rr|r}
}

\newcommand{\oneRowOpsTwo}[3] {
      #1 & #2 & #3 \\
}


\newcommand{\startRowOpsThree}{
  \left[
    \begin{array}{rrr|rrr}
}

\newcommand{\oneRowOpsThree}[6] {
      #1 & #2 & #3 & #4 & #5 & #6 \\
}


\newcommand{\startRowFour}{
  \left[
    \begin{array}{rrrr}
}

\newcommand{\oneRowFour}[4] {
      #1 & #2 & #3 & #4 \\
}

\newcommand{\stopRowOps}{
    \end{array}
  \right]
}




\begin{document}

\title{Ordinary Differential Equations}
\subtitle{Math 232 - Week 7, Day 3}

\author{Kelly Black}
\institute{Clarkson University}
\date{14 Oct 2011}

\begin{frame}
  \titlepage
\end{frame}

\begin{frame}
  \frametitle{Outline}
  \tableofcontents[pausesection,hideallsubsections]
\end{frame}


\section{Solutions to Second Order DEs}


\begin{frame}
  \frametitle{What is a solution to a general second order DE?}

  Solutions to
  \begin{eqnarray*}
    a y'' + by' + cy & = & 0.
  \end{eqnarray*}

  \uncover<2->
  {
    Assume
    \begin{eqnarray*}
      y & = & A e^{rt}, \\
      y' & = & r A e^{rt}, \\
      y'' & = & r^2 A e^{rt}.
    \end{eqnarray*}
    (If it works I am done!)
  }

\end{frame}


\begin{frame}
  \frametitle{Substitute back into the equation}

  The original de:
  \begin{eqnarray*}
    a y'' + by' + cy & = & 0, \\
    a r^2 A e^{rt} + b r A e^{rt} + c A e^{rt} & = & 0 \\
    \uncover<2->
    {
      A e^{rt} \left[ a r^2 + br + c \right] & = & 0.
    }
  \end{eqnarray*}

  \uncover<3->
  {
    This is true if
    \begin{eqnarray*}
      r & = & \frac{-b \pm \sqrt{b^2-4ac}}{2a}.
    \end{eqnarray*}
  }

  \uncover<4->
  {
    Three cases: distinct real, repeated roots, complex roots.
  }

\end{frame}

\section{Real, Distinct Roots}

\begin{frame}
  \frametitle{Case 1: Distinct, Real Roots}

  \begin{eqnarray*}
    y & = & C_1 e^{r_1 t} + C_2 e^{r_2 t}.
  \end{eqnarray*}

\end{frame}


\begin{frame}
  \frametitle{Example}

  Find the general solution to
  \begin{eqnarray*}
    y'' - 5y' + 6y & = & 0.
  \end{eqnarray*}

  \uncover<2->
  {
    \begin{eqnarray*}
      r^2 - 5r + 6 & = & 0, \\
      (r-3)(r-2) & = & 0.
    \end{eqnarray*}

    The general solution is
    \begin{eqnarray*}
      y & = & C_1 e^{3t} + C_2 e^{2t}
    \end{eqnarray*}

  }

\end{frame}


\begin{frame}
  \frametitle{Example}

  Find the general solution to
  \begin{eqnarray*}
    y'' - 2y' - 8y & = & 0.
  \end{eqnarray*}

  \uncover<2->
  {
    \begin{eqnarray*}
      r^2 - 2r - 8 & = & 0, \\
      (r-4)(r+2) & = & 0.
    \end{eqnarray*}

    The general solution is
    \begin{eqnarray*}
      y & = & C_1 e^{4t} + C_2 e^{-2t}
    \end{eqnarray*}

  }

\end{frame}


\section{Case 2: Repeated Roots}

\begin{frame}
  \frametitle{Case 2: Repeated Roots}

  \begin{eqnarray*}
    y & = & C_1 e^{rt} + C_2 t e^{rt}
  \end{eqnarray*}


\end{frame}


\begin{frame}
  \frametitle{Example}

  \begin{eqnarray*}
    y'' + 6y' + 9y & = & 0.
  \end{eqnarray*}

  \uncover<2->
  {
    \begin{eqnarray*}
      r^2 + 6r + 9 & = & 0, \\
      (r+3)^2 & = & 0.
    \end{eqnarray*}

    The general solution is
    \begin{eqnarray*}
      y & = & C_1 e^{-3t} + C_2 t e^{-3t}
    \end{eqnarray*}

  }

\end{frame}


\begin{frame}
  \frametitle{Nomenclature}

  For spring/mass systems (or RCL circuits)
  \begin{itemize}
  \item Distinct real, negative roots - \textit{Over damped}
  \item One repeated root - \textit{Critically damped}
  \item Complex roots   - \textit{Under damped}
  \end{itemize}

\end{frame}

\begin{frame}
  \frametitle{Example}

  \begin{eqnarray*}
    y'' + 4y' + 4y & = & 0, \\
    y(0) & = & 0, \\
    y'(0) & = & 3.
  \end{eqnarray*}

  \uncover<2->
  {
    \begin{eqnarray*}
      r^2 + 4r + 4 & = & 0, \\
      (r+2)^2 & = & 0.
    \end{eqnarray*}

    The general solution is
    \begin{eqnarray*}
      y & = & C_1 e^{-2t} + C_2 t e^{-2t}
    \end{eqnarray*}

  }

\end{frame}



\begin{frame}
  \frametitle{Solve for the constants}

  \begin{eqnarray*}
    y(0) & = & 0 \\
    & = & C_1 \\
    \Rightarrow C_1 & = & 0.
  \end{eqnarray*}

  \begin{eqnarray*}
    y(t) & = & C_2 t e^{-2t} \\
    y'(t) & = & C_2 e^{-2t} - 2 C_2 t e^{-2t}, \\
    y'(0) & = & C_2 \\
    \Rightarrow C_2 & = & 3.
  \end{eqnarray*}

  The solution is
  \begin{eqnarray*}
    y_h & = & 3 t e^{-2t}.
  \end{eqnarray*}

\end{frame}

\section{More Examples}

\begin{frame}
  \frametitle{Example}

  Find the general solution to
  \begin{eqnarray*}
    y'' + 10 y' + 24y & = & 0.
  \end{eqnarray*}

  \uncover<2->
  {
    \begin{eqnarray*}
      r^2 + 10r + 24 & = & 0, \\
      r & = & \frac{-10\pm\sqrt{100-4(24)}}{2} \\
      & = & -4,~ -6.
    \end{eqnarray*}

    The general solution is
    \begin{eqnarray*}
      y & = & C_1 e^{-4t} + C_2 e^{-6t}
    \end{eqnarray*}

  }
  

\end{frame}

\begin{frame}
  \frametitle{Linear Independence}

  Question: Are the functions linear independent?

  \uncover<2->
  {

    If we have $n$-derivatives we need $n$ linearly independent
    functions:
    \begin{eqnarray*}
      \mathrm{det}\arrayTwo{e^{-4t}}{e^{-6t}}{-4e^{-4t}}{-6e^{-6t}} 
      & = & -2 e^{-10t}.
    \end{eqnarray*}

  }
  
\end{frame}


\begin{frame}
  \frametitle{Example}

  Find the general solution to
  \begin{eqnarray*}
    y'' + y' - 2y & = & 0, \\
    y(0) & = & 1, \\
    y'(0) & = & 2.
  \end{eqnarray*}

  \uncover<2->
  {
    \begin{eqnarray*}
      r^2 + r - 2 & = & 0, \\
      (r-1)(r+2) & = & 0.
    \end{eqnarray*}

    The general solution is
    \begin{eqnarray*}
      y & = & C_1 e^{-2t} + C_2 e^{t}
    \end{eqnarray*}

  }
  

\end{frame}

\begin{frame}
  \frametitle{Find The Specific Solution}

  Apply the initial condition:
  \begin{eqnarray*}
    y(0) & = & C_1 + C_2 \\
    y'(0) & = & -2C_1 + C_2
  \end{eqnarray*}

  \begin{eqnarray*}
    \startRowOpsTwo
    \oneRowOpsTwo{1}{1}{1}
    \oneRowOpsTwo{-2}{1}{2}
    \stopRowOps \\
    \Rightarrow
    \startRowOpsTwo
    \oneRowOpsTwo{1}{0}{-\frac{1}{3}}
    \oneRowOpsTwo{0}{1}{\frac{4}{3}}
    \stopRowOps
  \end{eqnarray*}

  So the solution is
  \begin{eqnarray*}
    y & = & -\frac{1}{3} e^{-2t} + \frac{4}{3} e^{t}.
  \end{eqnarray*}
  

\end{frame}


\begin{frame}
  \frametitle{Higher Order Derivatives}

  No reason to stop at two derivatives! Find the general solution to 
  \begin{eqnarray*}
    y''' + 2 y'' - y' - 2y & = & 0.
  \end{eqnarray*}
  Assume that $y=Ae^{rt}$:
  \begin{eqnarray*}
    r^3 + 2 r^2 - r + 2 & = & 0, \\
    \uncover<2->
    {
      (r+2)(r-1)(r+1) & = & 0.
    }
  \end{eqnarray*}

  \uncover<3->
  {
    \begin{eqnarray*}
      y & = & C_1 e^{-2t} + C_2 e^{t} + C_3 e^{-t}.
    \end{eqnarray*}
  }
  
\end{frame}


\begin{frame}
  \frametitle{Higher Order Derivatives}

  Find the general solution to 
  \begin{eqnarray*}
    y''' + 11 y'' + 38 y' + 40y & = & 0.
  \end{eqnarray*}
  Assume that $y=Ae^{rt}$:
  \begin{eqnarray*}
    r^3 + 11 r^2 + 38 r + 40 & = & 0, \\
    \uncover<2->
    {
      (r+2)(r+4)(r+5) & = & 0.
    }
  \end{eqnarray*}

  \uncover<3->
  {
    \begin{eqnarray*}
      y & = & C_1 e^{-2t} + C_2 e^{-4t} + C_3 e^{-5t}.
    \end{eqnarray*}
  }
  
\end{frame}


\begin{frame}
  \frametitle{Higher Order Derivatives}

  Find the general solution to 
  \begin{eqnarray*}
    y''' + 3 y'' - 4y & = & 0.
  \end{eqnarray*}
  Assume that $y=Ae^{rt}$:
  \begin{eqnarray*}
    r^3 + 3 r^2 - 4 & = & 0, \\
    \uncover<2->
    {
      (r+2)^2(r-1) & = & 0.
    }
  \end{eqnarray*}

  \uncover<3->
  {
    \begin{eqnarray*}
      y & = & C_1 e^{-2t} + C_2 t e^{-2t} + C_3 e^{t}.
    \end{eqnarray*}
  }
  
\end{frame}


\begin{frame}
  \frametitle{Example}

  Find the general solution to
  \begin{eqnarray*}
    y'' + 7y' + 10y & = & 0, \\
    y(0) & = & 2, \\
    y'(0) & = & 3.
  \end{eqnarray*}

  \uncover<2->
  {
    \begin{eqnarray*}
      r^2 + 7r + 10 & = & 0, \\
      (r+2)(r+5) & = & 0.
    \end{eqnarray*}

    The general solution is
    \begin{eqnarray*}
      y & = & C_1 e^{-2t} + C_2 e^{-5t}
    \end{eqnarray*}

  }
  

\end{frame}


\begin{frame}
  \frametitle{Find The Specific Solution}

  Apply the initial condition:
  \begin{eqnarray*}
    y(0) & = & C_1 + C_2 \\
    y'(0) & = & -2C_1 - 5 C_2
  \end{eqnarray*}

  \begin{eqnarray*}
    \startRowOpsTwo
    \oneRowOpsTwo{1}{1}{2}
    \oneRowOpsTwo{-2}{-5}{3}
    \stopRowOps \\
    \Rightarrow
    \startRowOpsTwo
    \oneRowOpsTwo{1}{0}{\frac{13}{3}}
    \oneRowOpsTwo{0}{1}{-\frac{7}{3}}
    \stopRowOps
  \end{eqnarray*}

  So the solution is
  \begin{eqnarray*}
    y & = & \frac{13}{3} e^{-2t} - \frac{7}{3} e^{-5t}.
  \end{eqnarray*}
  

\end{frame}


\begin{frame}
  \frametitle{Example}

  Find the general solution to
  \begin{eqnarray*}
    2y'' + 7y' + 3y & = & 0, \\
    y(0) & = & 1, \\
    y'(0) & = & 4.
  \end{eqnarray*}

  \uncover<2->
  {
    \begin{eqnarray*}
      2r^2 + 7r + 3 & = & 0, \\
      r & = & \frac{-7\pm\sqrt{49-24}}{4}, \\
      r & = & -\half,~-3.
    \end{eqnarray*}

    The general solution is
    \begin{eqnarray*}
      y & = & C_1 e^{-t/2} + C_2 e^{-3t}
    \end{eqnarray*}

  }
  

\end{frame}


\begin{frame}
  \frametitle{Find The Specific Solution}

  Apply the initial condition:
  \begin{eqnarray*}
    y(0) & = & C_1 + C_2 \\
    y'(0) & = & -\half C_1 - 3 C_2
  \end{eqnarray*}

  \begin{eqnarray*}
    \startRowOpsTwo
    \oneRowOpsTwo{1}{1}{1}
    \oneRowOpsTwo{-\half}{-3}{4}
    \stopRowOps \\
    \Rightarrow
    \startRowOpsTwo
    \oneRowOpsTwo{1}{0}{\frac{14}{5}}
    \oneRowOpsTwo{0}{1}{-\frac{9}{5}}
    \stopRowOps
  \end{eqnarray*}

  So the solution is
  \begin{eqnarray*}
    y & = & \frac{14}{5} e^{-t/2} - \frac{9}{5} e^{-3t}.
  \end{eqnarray*}
  

\end{frame}



\end{document}

% LocalWords:  Clarkson pausesection hideallsubsections
